% Part: methods
% Chapter: induction
% Section: structural-induction

\documentclass[../../../include/open-logic-section]{subfiles}

\begin{document}

\olfileid{mth}{ind}{sti}

\olsection{结构归纳}

到目前为止，我们已经使用归纳法证明了关于所有自然数的结论。但存在一个与之类似的原理，可直接用于证明关于归纳定义集合中所有元素的结论。这种方法通常被称为\emph{结构}归纳，因为它依赖于归纳定义对象的\emph{结构}。

一般地，归纳定义由两部分给出：(a)~该集合的“初始”元素列表，以及 (b)~一系列从旧元素生成集合新元素的运算。例如，在好项（nice terms）的情形下，初始对象就是字母。我们只有一种运算，即 \begin{align*}
  o(s_1, s_2) = & [s_1 \circ s_2]
\end{align*}
你甚至可以把自然数集~$\Nat$ 本身看作是由归纳定义给出的：初始对象是~$0$，运算是后继函数~$x + 1$。

为了证明归纳定义集合中所有元素满足某条性质，即该集合中的每个元素都具有性质~$P$，我们必须：

\begin{enumerate}
\item 证明初始对象具有性质~$P$；
\item 证明对于每个运算~$o$，如果其自变元具有性质~$P$，那么运算结果也具有性质~$P$。
\end{enumerate}

例如，为了证明关于所有好项的某条性质，我们会证明它对所有字母成立，并且只要它对 $s_1$ 和 $s_2$ 分别成立，那么它对 $[s_1 \circ s_2]$ 也成立。

\begin{prop}
  在任何好项~$t$ 中，$[$ 的数量等于 $]$ 的数量。
\end{prop}

\begin{proof}
我们使用结构归纳法。好项是归纳定义的，其中字母是初始对象，运算 $o$ 用于从旧的好项构造新的好项。
\begin{enumerate}
\item 该断言对每个字母都成立，因为单个字母中 $[$ 的数量是~$0$，$]$ 的数量也是~$0$。
\item 假设 $s_1$ 中 $[$ 的数量等于 $]$ 的数量，且 $s_2$ 也是如此。$o(s_1, s_2)$，即 $[s_1 \circ s_2]$ 中 $[$ 的数量，是 $s_1$ 和 $s_2$ 中 $[$ 的数量之和再加一。$o(s_1, s_2)$ 中 $]$ 的数量是 $s_1$ 和 $s_2$ 中 $]$ 的数量之和再加一。因此，$o(s_1, s_2)$ 中 $[$ 的数量等于 $o(s_1,s_2)$ 中 $]$ 的数量。
\end{enumerate}
\end{proof}

\begin{prob}
  用结构归纳法证明：没有好项以~$]$ 开头。
\end{prob}

让我们再给出一个结构归纳法的证明：符号串~$t$ 的\emph{真初始段}是指任何从左向右逐符号与 $t$ 相同，但 $t$ 更长的符号串~$s$。例如，$[a \circ {}$ 是 $[a \circ b]$ 的一个真初始段，但 $[b \circ {}$（它们在第二个符号处不同）和 $[a \circ b]$（它们长度相同）都不是。

\begin{prop}\ollabel{prop:initial}
  好项~$t$ 的每个真初始段中，$[$ 的数量多于 $]$ 的数量。
\end{prop}

\begin{proof}
  对~$t$ 进行归纳：
  \begin{enumerate}
  \item $t$ 本身是一个字母：那么 $t$ 没有真初始段。
  \item $t = [s_1 \circ s_2]$，其中 $s_1$ 和 $s_2$ 是某些好项。如果 $r$ 是 $t$ 的一个真初始段，则有几种可能情况：
    \begin{enumerate}
    \item $r$ 就是 $[$：那么 $r$ 中 $[$ 的数量比 $]$ 多一个。
    \item $r$ 是 $[r_1$，其中 $r_1$ 是~$s_1$ 的一个真初始段：由于 $s_1$ 是好项，根据归纳假设，$r_1$ 中 $[$ 多于 $]$，$[r_1$ 亦然。
    \item $r$ 是 $[s_1$ 或 $[s_1 \circ {}$：根据前一个结果，$s_1$ 中 $[$ 和 $]$ 的数量相等；所以 $[s_1$ 或 $[s_1 \circ {}$ 中 $[$ 的数量比 $]$ 的数量多一个。
    \item $r$ 是 $[s_1 \circ r_2$，其中 $r_2$ 是~$s_2$ 的一个真初始段：根据归纳假设，$r_2$ 包含的 $[$ 多于 $]$。根据前一个结果，$s_1$ 中 $[$ 和 $]$ 的数量相等。所以 $[s_1 \circ r_2$ 中 $[$ 的数量大于 $]$ 的数量。
    \item $r$ 是 $[s_1 \circ s_2$：根据前一个结果，$s_1$ 中 $[$ 和 $]$ 的数量相等，$s_2$ 也是如此。所以 $[s_1 \circ s_2$ 中 $[$ 的数量比 $]$ 多一个。
    \end{enumerate}
  \end{enumerate}
\end{proof}

\end{document}